% Section owner: Puneet (protocol) and Rohan (evaluation metrics)
\section{Protocol and Evaluation}
\label{sec:p2protocol}
\textbf{Selection.} All choices use validation data only. The final deep model is the configuration with the
highest mean validation macro-F1 over the three seeds (ties: accuracy, then fewer parameters); this rule is
implemented in code and was fixed before the runs.

\textbf{Freeze and test gate.} After selection and hybrid tuning, one file records the selected configuration,
the hybrid hyper-parameters, $\tau$, both recipes, the split hash and the list of \NGateModels{} models allowed on
the test set: for each seed the re-trained Phase 1 CNN, the final deep model, H1, H2 and H3, and H0. The test
set is loaded inside the gate, which refuses to run without the freeze file, with an existing flag file, or when
the evaluated models differ from the frozen list. The gate ran once on all 10{,}000 test images.

\textbf{Test-set reuse.} The same test images were evaluated once at the Phase 1 gate, and the Phase 1 error
analysis motivated Phase 2. Phase 2 test numbers are therefore slightly optimistic as estimates for new data.
Comparisons between Phase 2 models remain fair because all of them were selected on validation only.

\textbf{Two kinds of uncertainty.} Training randomness is reported as mean $\pm$ sample standard deviation over
three seeds. Test-set sampling is reported as a 95\% bootstrap interval (1{,}000 resamples) for seed 42. Pairs of
models are compared with McNemar's test on seed 42, Bonferroni-corrected for \NComparisons{} planned comparisons.
With three seeds the standard deviation is itself uncertain, so we use it descriptively.

\textbf{Further measurements.} Per-class recall; in-group confusion, the share of T-shirt, Pullover, Coat and
Shirt images assigned to another class of that group; the expected calibration error with 15 confidence
bins~\cite{guo2017calibration}; accuracy under five corruptions on a stratified subset of 2{,}000 test images; and
parameter counts and inference time.

\textbf{Compute.} All \NRuns{} runs were trained on one Tesla T4 GPU in \TrainMinutesTotal{} minutes of training
time in total (about \TrainSecondsDeeper{}~s per run for the final architecture and \TrainSecondsPone{}~s for the
Phase 1 CNN). Runs are stored individually and can be resumed.
